%=====================================================================
\section{The borehole resistance network}
\label{sec:network}

At a given axial station the heat leaving the secondary fluid crosses, in
series: the internal convective film, the fouling layer, the tube wall, and the
melt layer standing on a finned outer surface. All four are referred to the
inner tube area, so that a single coefficient $\Ui$ drives the segment march of
\S\ref{sec:front:march}.

%---------------------------------------------------------------------
\subsection{Internal convection}

With $D_{i}=2r_{i}$ and $\md$ the flow in one tube,
\begin{equation}
  \mathrm{Re} \;=\; \frac{4\md}{\pi D_{i}\,\mu},
  \qquad
  \mathrm{Pr} \;=\; \frac{c_{p}\mu}{k}.
\end{equation}
The Nusselt number is the Gnielinski correlation in turbulent flow and a fixed
fully-developed value in laminar flow:
\begin{equation}
  \mathrm{Nu} =
  \begin{dcases}
    \mathrm{Nu}_{\mathrm{lam}}, & \mathrm{Re} < 2300, \\[4pt]
    \frac{(f_{G}/8)(\mathrm{Re}-1000)\,\mathrm{Pr}}
         {1 + 12.7\sqrt{f_{G}/8}\,\bigl(\mathrm{Pr}^{2/3}-1\bigr)},
      & \mathrm{Re} \geq 2300,
  \end{dcases}
  \qquad
  f_{G} = \bigl(0.79\ln\mathrm{Re} - 1.64\bigr)^{-2},
  \label{eq:gnielinski}
\end{equation}
and
\begin{equation}
  h_{i} \;=\; \frac{\mathrm{Nu}\,k}{D_{i}} .
\end{equation}

The laminar constant differs between the two implementations:
$\mathrm{Nu}_{\mathrm{lam}}=4.36$ in \code{well\_marched} (uniform wall flux) and
$3.66$ in the legacy \code{temperature\_profile\_melt} (uniform wall
temperature). For a melting boundary the wall is closer to isothermal, so $3.66$
is the better constant of the two; the discrepancy is recorded in
\S\ref{sec:defects} and is inconsequential at the design point, where the flow
is strongly turbulent.

%---------------------------------------------------------------------
\subsection{Wall, fouling and melt layer}

The wall and fouling resistances, referred to the inner area, are
\begin{equation}
  R'_{\mathrm{wall}} \;=\; \frac{r_{i}}{k_{w}}\,\ln\!\frac{r_{e}}{r_{i}},
  \qquad
  R'_{f} \;=\; R''_{f,i},
  \qquad
  R'_{\mathrm{conv}} \;=\; \frac{1}{h_{i}} .
  \label{eq:R-wall}
\end{equation}

The melt layer is treated as a steady cylindrical conduction shell between
$r_{e}$ and $r_{e}+\delta$. Expressed as a conductance on the tube outer
surface,
\begin{equation}
  h_{e} \;=\; \frac{\km}{r_{e}\,\ln\!\bigl(1+\delta/r_{e}\bigr)} .
  \label{eq:he}
\end{equation}
As $\delta\to 0$ the layer vanishes and $h_{e}\to\infty$; the code caps this at
$10^{9}$, which is the correct limit and not, as an earlier audit asserted, a
defect. What \emph{is} defective is that error paths also route into that
branch, so a failed solve is reported as maximum heat transfer
(\S\ref{sec:defects}).

Equation~\eqref{eq:he} carries an assumption that becomes important in
\S\ref{sec:defects:h3}: it presumes the \emph{entire} azimuth
$2\pi(r_{e}+\delta)$ is available for conduction at every radius. That holds
only while the melt front is free of its neighbours.

%---------------------------------------------------------------------
\subsection{Which shell does the heat cross?}
\label{sec:network:shell}

Equation~\eqref{eq:he} is a shell between $r_{e}$ and $r_{e}+\delta$, so it
poses a question that is easy to pass over: \emph{which material is that shell
made of, and how thick is it?} The answer is not the same in the two
half-cycles, and getting it wrong is invisible in every melting result.

The front always grows \emph{away from the tube}, because the tube is the driven
boundary in both directions.

\begin{center}
\begin{tabular}{@{}l l l l@{}}
\toprule
 & shell against the tube & thickness from & conductivity \\
\midrule
melting  & liquid        & $\Amelt$               & $k_{l}$ \\
freezing & frozen solid  & $A_{\mathrm{avail}}-\Amelt$ & $k_{s}$ \\
\bottomrule
\end{tabular}
\end{center}

\noindent
While melting, heat crosses the melt to reach the remaining solid --- the
familiar picture. While freezing, solidification also nucleates at the wall, so
heat crosses the \emph{frozen} shell, and the liquid is displaced outward beyond
the front where it plays no part in the conduction path at all.

Written this way the four limiting states need no special handling:

\begin{center}
\begin{tabular}{@{}l l l l@{}}
\toprule
state & direction & $\delta$ & PCM-side resistance \\
\midrule
fully solid  & melting  & $0$                    & none; the front is at the tube \\
fully melted & melting  & $\delta_{\mathrm{merge}}$ & the full liquid annulus \\
fully melted & freezing & $0$                    & none; nothing has frozen \\
fully solid  & freezing & $\delta_{\mathrm{merge}}$ & the full frozen annulus \\
\bottomrule
\end{tabular}
\end{center}

\noindent
Function: \code{conduction\_shell}. The direction is taken from the sign of
$T_{j}-\Tpcm$, which is known before the conductance is formed and has the same
sign as $q'$.

\subsubsection*{What this replaced, and what it was worth}

Until v0.7 the \emph{melted} thickness was used in both directions, with
$\km=k_{l}$ whenever any melt existed. On charging that is correct. On
discharging it puts the conduction path on the wrong side of the front, and the
two thicknesses are complementary, so the error is a \emph{time reversal} rather
than a scale factor:

\begin{center}
\begin{tabular}{@{}r r r r r r@{}}
\toprule
$\varepsilon_{\mathrm{local}}$ & $\delta$ used & $\delta$ wanted
 & $R'$ used & $R'$ wanted & ratio \\
 & \si{\milli\meter} & \si{\milli\meter}
 & \multicolumn{2}{c}{\si{\meter\kelvin\per\watt}} & \\
\midrule
1.00 & 23.37 &  0.00 & 0.2639 & 0.0000 & $\infty$ \\
0.50 & 14.32 & 14.32 & 0.1833 & 0.1375 & 1.33 \\
0.08 &  3.48 & 22.01 & 0.0541 & 0.1897 & 0.28 \\
0.00 &  0.00 & 23.37 & 0.0000 & 0.1979 & 0 \\
\bottomrule
\end{tabular}
\end{center}

\noindent
A discharge runs down that table. The modelled resistance \emph{fell} by a
factor of five as the store froze, where it should rise from nothing to its
maximum; a segment frozen solid was given \textbf{zero} PCM-side resistance at
precisely the point where the physical resistance is largest.

The consequence is visible in $\Ui$ and was for a long time mistaken for
physics. Mid-well at the end of each half-cycle the discharge value exceeded the
charge value by a factor of \textbf{seven}, and \S\ref{sec:defects} recorded
that as the melt layer refreezing --- solid PCM conducting better than liquid.
But $k_{s}/k_{l}=1.33$, and nothing accounted for the remaining $5.3$. With the
shell taken on the correct side, the two half-cycles become near mirror images:
at $t\to0$ the values agree to \SI{0.5}{\percent}, because the shell has zero
thickness at the start of \emph{both}; at the end the ratio is $1.265$, which is
$k_{s}/k_{l}$ less the fact that the frozen shell does not grow quite as thick
as the melted one.

\begin{center}
\begin{tabular}{@{}l r r@{}}
\toprule
at the design point & melt side & directional \\
\midrule
$\Ui$ mid-well, $t\to0$, charge / discharge & 822 / 117 & 823 / 829 \\
$\Ui$ mid-well, end, charge / discharge     & 117 / 821 & 117 / 148 \\
CSS deviation                                & \SI{+1.23}{\percent} & \SI{+4.73}{\percent} \\
residual $\varepsilon$ at end of discharge   & 0.0823 & 0.0094 \\
$\varepsilon$ at end of charge               & 1.0000 & 0.9706 \\
$\Nwells$                                    & 11.743 & 11.626 \\
\bottomrule
\end{tabular}
\end{center}

\noindent
The direction of that change is worth dwelling on, because it is the opposite of
what one would guess. The retired treatment was too \emph{optimistic} at the end
of the discharge, which is where the rate limit bites --- so the correction
ought to reduce the delivered energy. It increases it. At the \emph{start} of
the discharge the retired treatment was infinitely too pessimistic, and that is
when the driving difference is largest and most of the energy moves; the early
over-penalty dominates the integral. The store also now very nearly refreezes,
so the discharge is closer to inventory-limited at the design point than it
appeared, although the rate-limit diagnostic survives: the residual
$\varepsilon$ still rises with well count.

\code{Case.front\_geometry = 'melt\_side'} reproduces the retired behaviour
exactly, which is how this change was separated from the closure change of
\S\ref{sec:cycles:closure} in the frozen fixture.

\subsubsection*{And when there is no front at all}

The rule above answers ``which side of the front'', which presupposes a front. A
cell that is \emph{entirely} one phase --- all subcooled solid, or all
superheated liquid --- has none, and asking the rule anyway returns a resistance
that depends on the \emph{direction} of the heat flow: zero one way, the full
annulus the other. Same body, same material, two answers.

What changes at a branch boundary is not the geometry but the \emph{meaning of
the state variable}. In two-phase, $\Tpcm=\Tm$ is the temperature of the
\emph{front}, and the resistance that belongs with it is tube-to-front. In
single phase there is no front: $\Tpcm$ is the \emph{volume mean} of the cell,
and the resistance that belongs with a mean is mean-to-surface.

The derivation is elementary but it is set out in full below, because the
result is used everywhere in the single-phase branches and because each step
carries a modelling choice that is easy to make silently.

\subsubsection*{Step 1: what is being computed}

The quantity wanted is a \emph{mean-to-surface} resistance,
\begin{equation}
  R'_{\mathrm{bulk}} \;\equiv\; \frac{\overline{T}-T_{s}}{Q'},
  \qquad
  \overline{T} \;\equiv\; \frac{1}{V}\int_{V} T\,\mathrm{d}V,
  \label{eq:bulk-def}
\end{equation}
and not the surface-to-surface resistance of Eq.~\eqref{eq:he}. The averaging
is not a choice of convenience: the state variable is $E'$, the \emph{total}
energy of the cell per unit tube length, so the temperature the formulation
inverts out of it is by construction the volume mean. Any other weighting would
be inconsistent with the definition of $E'$ itself.

\subsubsection*{Step 2: a field to average}

A mean needs a profile, and a profile needs a heat-flow model. Take the cell as
an annulus $r_{e}\le r\le r_{o}$, adiabatic at $r_{o}$ --- the neighbouring cell
is identical, so that boundary is a symmetry plane --- and quasi-steady, so that
its stored energy changes at a uniform volumetric rate
$s=\rho c\,\mathrm{d}T/\mathrm{d}t$ while the tube carries the whole flux. A
uniformly distributed storage rate is mathematically a uniform volumetric
source, and the problem becomes the classical one of steady conduction with
uniform generation:
\begin{equation}
  \frac{k}{r}\frac{\mathrm{d}}{\mathrm{d}r}
    \left(r\frac{\mathrm{d}T}{\mathrm{d}r}\right) + s \;=\; 0,
  \qquad
  \left.\frac{\mathrm{d}T}{\mathrm{d}r}\right|_{r_{o}} = 0 .
  \label{eq:bulk-pde}
\end{equation}

\subsubsection*{Step 3: integrate once --- the flux}

One integration of Eq.~\eqref{eq:bulk-pde} gives
$rk\,\mathrm{d}T/\mathrm{d}r=-sr^{2}/2+C_{1}$, and the adiabatic condition at
$r_{o}$ fixes $C_{1}=sr_{o}^{2}/2$:
\begin{equation}
  -k\,2\pi r\,\frac{\mathrm{d}T}{\mathrm{d}r}
  \;=\; -\,s\,\pi\left(r_{o}^{2}-r^{2}\right).
  \label{eq:bulk-ode}
\end{equation}
Equation~\eqref{eq:bulk-ode} states that whatever crosses radius $r$ is what the
material beyond it stores. Evaluated at $r=r_{e}$ it returns the cell energy
balance, $\left|Q'\right|=s\pi(r_{o}^{2}-r_{e}^{2})$, which is a check that $s$
and $Q'$ are the same statement rather than two independent quantities.

\subsubsection*{Step 4: integrate again --- the profile}

Writing $\theta(r)=T(r)-T(r_{e})$ and integrating Eq.~\eqref{eq:bulk-ode} from
the tube outward,
\begin{equation}
  \theta(r) \;=\; \frac{s}{2k}\int_{r_{e}}^{r}
    \left(\frac{r_{o}^{2}}{r'}-r'\right)\mathrm{d}r'
  \;=\; \frac{s}{2k}
  \left[r_{o}^{2}\ln\frac{r}{r_{e}}-\frac{r^{2}-r_{e}^{2}}{2}\right].
  \label{eq:bulk-profile}
\end{equation}

\subsubsection*{Step 5: average over the volume}

With the annular element $\mathrm{d}V=2\pi r\,\mathrm{d}r$ per unit length,
\begin{equation}
  \overline{\theta} \;=\;
  \frac{\displaystyle\int_{r_{e}}^{r_{o}}\theta(r)\,2\pi r\,\mathrm{d}r}
       {\pi\left(r_{o}^{2}-r_{e}^{2}\right)}
  \;=\; \frac{2}{r_{o}^{2}-r_{e}^{2}}
        \int_{r_{e}}^{r_{o}}\theta(r)\,r\,\mathrm{d}r ,
  \label{eq:bulk-average}
\end{equation}
which requires only two elementary integrals, with $\beta=r_{o}/r_{e}$:
\begin{equation}
  \int_{r_{e}}^{r_{o}}\! r\ln\frac{r}{r_{e}}\,\mathrm{d}r
    = \frac{r_{o}^{2}}{2}\ln\beta-\frac{r_{o}^{2}-r_{e}^{2}}{4},
  \qquad
  \int_{r_{e}}^{r_{o}}\!\left(r^{2}-r_{e}^{2}\right) r\,\mathrm{d}r
    = \frac{r_{o}^{4}-r_{e}^{4}}{4}
      -\frac{r_{e}^{2}\left(r_{o}^{2}-r_{e}^{2}\right)}{2}.
  \label{eq:bulk-integrals}
\end{equation}

\subsubsection*{Step 6: form the resistance and non-dimensionalise}

Dividing Eq.~\eqref{eq:bulk-average} by $Q'=s\pi(r_{o}^{2}-r_{e}^{2})$ removes
$s$ entirely --- as it must, or the ``resistance'' would depend on how hard the
cell was being driven. Writing $R'_{\mathrm{bulk}}=S/2\pi k$ defines the shape
factor $S=2\pi k R'_{\mathrm{bulk}}$, and substituting $r_{o}=\beta r_{e}$ into
Eqs.~\eqref{eq:bulk-average}--\eqref{eq:bulk-integrals} gives
\begin{equation}
  \boxed{\;S \;=\;
  \frac{4\beta^{4}\ln\beta-3\beta^{4}+4\beta^{2}-1}
       {4\left(\beta^{2}-1\right)^{2}}\;}
  \label{eq:shape-factor}
\end{equation}

\noindent
$S$ depends on \textbf{geometry alone} --- $k$ cancels --- so one shape factor
serves both phases and only the conductivity changes between them.

\subsubsection*{Step 7: two checks}

\noindent\textbf{(a) The thin-annulus limit.} Write the wall thickness of the
annulus as $L=r_{o}-r_{e}$, so that $\beta=1+L/r_{e}$, and expand
Eq.~\eqref{eq:shape-factor} for $L/r_{e}\ll1$. This gives $S\to L/3r_{e}$ and
therefore
\begin{equation}
  R'_{\mathrm{bulk}} \;\longrightarrow\;
  \frac{L/3r_{e}}{2\pi k}
  \;=\; \frac{L}{3k\left(2\pi r_{e}\right)}
  \;=\; \frac{L}{3kA},
  \label{eq:slab-limit}
\end{equation}
with $A$ the surface area per unit length: the mean-to-surface resistance of a
plane slab with uniform generation. The convergence is fast --- at $L/r_{e}=0.2$
the ratio $S/(L/3r_{e})$ is already $0.9965$, and at $0.05$ it is $0.9998$. The
cylindrical result is the familiar $1/3$ deformed by curvature, and nothing
more.

\medskip
\noindent\textbf{(b) Against the surface-to-surface annulus.} For the default
cell $\beta=2.1086$ and $S=0.34672$, against $\ln\beta=0.74604$ in
Eq.~\eqref{eq:he}: the mean-to-surface resistance is $46.5\%$ of the
surface-to-surface one. Smaller is the only defensible sign, and not by a
factor so near unity that the distinction would be cosmetic. In absolute terms,

\begin{center}
\begin{tabular}{@{}l r r@{}}
\toprule
 & $R'_{\mathrm{bulk}}=S/2\pi k$ & $R'_{\mathrm{annulus}}=\ln\beta/2\pi k$ \\
 & [\si{\meter\kelvin\per\watt}] & [\si{\meter\kelvin\per\watt}] \\
\midrule
solid, $k_{s}=\SI{0.60}{\watt\per\meter\per\kelvin}$  & 0.0920 & 0.1979 \\
liquid, $k_{l}=\SI{0.45}{\watt\per\meter\per\kelvin}$ & 0.1226 & 0.2639 \\
\bottomrule
\end{tabular}
\end{center}

\noindent
Equation~\eqref{eq:shape-factor} was additionally checked against numerical
quadrature of Eq.~\eqref{eq:bulk-average} to six decimals, and every step above
against a symbolic algebra system.

\subsubsection*{Step 8: back into the network as a thickness}

The network of \S\ref{sec:network:fins} takes a thickness rather than a
resistance, and wraps it in the fin efficiency and the finned-area ratio. Fed
back as an equivalent thickness,
\begin{equation}
  \ln\left(1+\frac{\delta_{\mathrm{eq}}}{r_{e}}\right)=S
  \qquad\Longrightarrow\qquad
  \delta_{\mathrm{eq}} = r_{e}\left(e^{S}-1\right) = \SI{8.74}{\milli\meter},
  \label{eq:delta-eq}
\end{equation}
against $\delta_{\mathrm{merge}}=r_{o}-r_{e}=\SI{23.37}{\milli\meter}$. The
single-phase branches then stay inside exactly the same machinery as the
two-phase ones instead of needing a second path. Like $S$, $\delta_{\mathrm{eq}}$
is a property of the geometry alone. Functions: \code{bulk\_shape\_factor},
\code{bulk\_equivalent\_delta}.

\begin{center}
\begin{tabular}{@{}l l l@{}}
\toprule
state & before v0.8 & v0.8 \\
\midrule
subcooled solid, heated    & $\delta=0$, no resistance &
  $\delta_{\mathrm{eq}}$, $k_{s}$ \\
subcooled solid, cooled    & $\delta_{\mathrm{merge}}$, the full annulus &
  $\delta_{\mathrm{eq}}$, $k_{s}$ \\
superheated liquid, cooled & $\delta=0$, no resistance &
  $\delta_{\mathrm{eq}}$, $k_{l}$ \\
superheated liquid, heated & $\delta_{\mathrm{merge}}$, the full annulus &
  $\delta_{\mathrm{eq}}$, $k_{l}$ \\
\bottomrule
\end{tabular}
\end{center}

\noindent
The two retired values \emph{bracket} the answer from opposite sides --- zero,
and $2.152$ times too large --- and which one a cell received depended on the
direction of the heat flow.

\medskip
\noindent\textbf{The jump at a branch boundary is real, and it is not an
artefact.} As $\varepsilon_{\mathrm{local}}\to1^{-}$ while melting, the front is
at $r_{o}$ and $R'=\ln\beta/2\pi k_{l}$; the instant the cell becomes
superheated, $R'$ drops to $S/2\pi k_{l}$. Nothing physical has moved. What
changed is that $\Tpcm$ stopped being the front temperature and became the
volume mean, and the resistance follows its reference. The same occurs at
$\varepsilon_{\mathrm{local}}\to0$, where the mechanism switches from
\emph{warming a body} to \emph{moving a front}. The flux never becomes
unbounded, because Eq.~\eqref{eq:K-derived} caps it at the capacity rate.

\medskip
\noindent
The effect on the integrated result is modest --- the CSS deviation moves
\SI{+4.73}{\percent} to \SI{+4.53}{\percent}, because the two errors partly
cancel over a half-cycle. The qualitative effect is larger, and is set out in
\S\ref{sec:defects}: the spurious infinite conductance at
$\varepsilon_{\mathrm{local}}=0$ and $1$ is gone, and with it the claim that the
exchanger is at its best at $t\to0$. It is at its best shortly \emph{after}
melting begins, when a segment has a front at the tube and almost nothing
between them.

\subsubsection*{The honest limit}

One thickness describes one front. After the first half-cycle from a fully solid
store the geometry is generally \emph{three-region}: at cyclic steady state the
charge begins with residual liquid in the outer part of the cell, so melting
produces liquid at the tube, solid in the middle and liquid outside. No single
$\delta$ represents that, and neither treatment can. The arrangement above is
the better approximation, not a correct one --- which is what hypothesis H5
actually assumes, stated plainly.

%---------------------------------------------------------------------
\subsection{Fin efficiency and the finned surface}
\label{sec:network:fins}

The fins are straight longitudinal fins of rectangular profile, treated with the
adiabatic-tip correction $L_{c}=L_{f}+t_{f}/2$:
\begin{equation}
  m \;=\; \sqrt{\frac{2h_{e}}{k_{w}\,t_{f}}},
  \qquad
  \etaf \;=\; \frac{\tanh(m L_{c})}{m L_{c}} .
  \label{eq:fin-eff}
\end{equation}

With the fin surface and total outer surface per tube
\begin{equation}
  A_{f} \;=\; L_{\mathrm{tube}}\,\bigl(2L_{f}+t_{f}\bigr),
  \qquad
  A_{t} \;=\; 2L_{\mathrm{tube}}\,\bigl(\pi r_{e} + n_{f} L_{f}\bigr),
  \label{eq:areas}
\end{equation}
the overall surface efficiency is
\begin{equation}
  \etao \;=\; 1 - \frac{n_{f}A_{f}}{A_{t}}\,\bigl(1-\etaf\bigr),
  \label{eq:eta-o}
\end{equation}
and the wetted perimeter of the finned tube is
\begin{equation}
  P_{T} \;=\; 2\pi r_{e} + 2 n_{f} L_{f} .
  \label{eq:perim}
\end{equation}

Equation~\eqref{eq:perim} is the quantity at the centre of the v0.2 correction.
At the design geometry $P_{T}=\SI{0.4924}{\meter}$, against a bare-tube
perimeter $2\pi r_{e}=\SI{0.1324}{\meter}$ --- a factor of $3.72$. The legacy
melt-front model of \S\ref{sec:front:legacy} uses the bare perimeter while the
network above uses $P_{T}$, and nothing in v0.1 compared them.

\medskip
\noindent\textbf{How the fins enter, and how they do not.} Equations
\eqref{eq:fin-eff}--\eqref{eq:perim} represent a fin as \emph{additional surface
at the tube}: more perimeter, discounted by an efficiency. That is the whole of
the fin model. A fin is \emph{not} represented as a radial conduction path
carrying heat outward to PCM the tube cannot easily reach. The distinction is
immaterial while the melt layer is thin and becomes material once the fronts
approach one another, which is the situation analysed in
\S\ref{sec:defects:h3}.

Note also that $k_{w}$ appears twice with different physical meanings: as the
tube-wall conductivity in Eq.~\eqref{eq:R-wall} and as the \emph{fin}
conductivity in Eq.~\eqref{eq:fin-eff}. A single argument carries both. This is
what allowed the PCM liquid conductivity to be passed in that slot for months
without contradiction (\S\ref{sec:defects:kw}).

%---------------------------------------------------------------------
\subsection{Overall coefficient}

The melt-layer resistance referred to the inner area carries the ratio of
perimeters and the surface efficiency,
\begin{equation}
  R'_{\mathrm{outer}} \;=\; \frac{2\pi r_{i}}{P_{T}\,\etao\,h_{e}},
  \label{eq:R-outer}
\end{equation}
and the overall coefficient is the inverse of the series sum,
\begin{equation}
  \boxed{\;
  \Ui \;=\;
  \Bigl[\,
    \underbrace{\tfrac{1}{h_{i}}}_{\text{convection}}
    +\;\underbrace{R''_{f,i}}_{\text{fouling}}
    +\;\underbrace{\tfrac{r_{i}}{k_{w}}\ln\tfrac{r_{e}}{r_{i}}}_{\text{wall}}
    +\;\underbrace{\tfrac{2\pi r_{i}}{P_{T}\etao h_{e}}}_{\text{melt layer}}
  \,\Bigr]^{-1} .
  \;}
  \label{eq:Ui}
\end{equation}

\emph{Source:} \code{\_legacy\_physics.compute\_U\_i}.

%---------------------------------------------------------------------
\subsection{Outer-surface temperature}
\label{sec:network:Tre}

The legacy melt-front model is driven by the tube outer-surface temperature
rather than by the heat flow. Splitting the network at $r_{e}$, with
\begin{equation}
  R_{1} \;=\; R''_{f,i} + \frac{1}{h_{i}}
              + \frac{r_{i}}{k_{w}}\ln\frac{r_{e}}{r_{i}},
  \qquad
  R_{2} \;=\; \frac{2\pi r_{i}}{P_{T}\,h_{e}\,\etao},
\end{equation}
and $T_{\mathrm{avg}}=\tfrac{1}{2}(T_{j}+T_{j+1})$ the mean fluid temperature
over the segment, a series voltage-divider gives
\begin{equation}
  \Tre \;=\; T_{\mathrm{avg}}
             - \frac{R_{1}\bigl(T_{\mathrm{avg}}-\Tm\bigr)}{R_{1}+R_{2}} .
  \label{eq:Tre}
\end{equation}

\emph{Source:} \code{\_legacy\_physics.compute\_T\_re}.

Equations~\eqref{eq:Ui} and~\eqref{eq:Tre} are solved together with the melt
front in a fixed-point loop at each segment, since $h_{e}$ depends on $\delta$,
$\delta$ depends on $\Tre$, and $\Tre$ depends on $h_{e}$. The loop runs to a
tolerance $\delta_{\mathrm{tol}}$ or a maximum iteration count, warm-started
from the previous segment.

Equation~\eqref{eq:Tre} is required \emph{only} by the legacy front. The
formulations of \S\ref{sec:front:marched} and \S\ref{sec:front:enthalpy} advance
the front from $q'$ directly and never form $\Tre$, which removes the
fixed-point loop along with it.
